\documentclass[12pt, oneside]{article}

\usepackage[letterpaper, scale=0.89, centering]{geometry}
\usepackage{fancyhdr}
\setlength{\parindent}{0em}
\setlength{\parskip}{1em}

\pagestyle{fancy}
\fancyhf{}
\renewcommand{\headrulewidth}{0pt}
\rfoot{\href{https://creativecommons.org/licenses/by-nc-sa/2.0/}{CC BY-NC-SA 2.0} Version \today~(\thepage)}

\usepackage{amssymb,amsmath,pifont,amsfonts,comment,enumerate,enumitem}
\usepackage{currfile,xstring,hyperref,tabularx,graphicx,wasysym}
\usepackage[labelformat=empty]{caption}
\usepackage{xcolor}
\usepackage{multicol,multirow,array,listings,tabularx,lastpage,textcomp,booktabs}

\lstnewenvironment{algorithm}[1][] {   
    \lstset{ mathescape=true,
        frame=tB,
        numbers=left, 
        numberstyle=\tiny,
        basicstyle=\rmfamily\scriptsize, 
        keywordstyle=\color{black}\bfseries,
        keywords={,procedure, div, for, to, input, output, return, datatype, function, in, if, else, foreach, while, begin, end, }
        numbers=left,
        xleftmargin=.04\textwidth,
        #1
    }
}
{}
\lstnewenvironment{java}[1][]
{   
    \lstset{
        language=java,
        mathescape=true,
        frame=tB,
        numbers=left, 
        numberstyle=\tiny,
        basicstyle=\ttfamily\scriptsize, 
        keywordstyle=\color{black}\bfseries,
        keywords={, int, double, for, return, if, else, while, }
        numbers=left,
        xleftmargin=.04\textwidth,
        #1
    }
}
{}

\newcommand\abs[1]{\lvert~#1~\rvert}
\newcommand{\st}{\mid}

\newcommand{\A}[0]{\texttt{A}}
\newcommand{\C}[0]{\texttt{C}}
\newcommand{\G}[0]{\texttt{G}}
\newcommand{\U}[0]{\texttt{U}}

\newcommand{\cmark}{\ding{51}}
\newcommand{\xmark}{\ding{55}}

 
\begin{document}
\begin{flushright}
    \StrBefore{\currfilename}{.}
\end{flushright} 
\subsection*{Week 2 at a glance}

\subsubsection*{We will be learning and practicing to:}
\begin{itemize}
\item Model systems with tools from discrete mathematics and reason about implications of modelling choices. Explore applications in CS through multiple perspectives, including software, hardware, and theory.
\begin{itemize}
   \item Selecting and representing appropriate data types and using notation conventions to clearly communicate choices
   \item Determining the properties of positional number representations, including overflow and bit operations
\end{itemize}

\item Translate between different representations to illustrate a concept.

\begin{itemize}
   \item Translating between symbolic and English versions of statements using precise mathematical language
   \item Tracing algorithms specified in pseudocode
   \item Representing numbers using positional representations, including decimal, binary, hexadecimal, fixed-width representations, and 2s complement
\end{itemize}


\item Use precise notation to encode meaning and present arguments concisely and clearly
\begin{itemize}
    \item Precisely describing a set using appropriate notation e.g. roster method, set builder notation, and recursive definitions
    \item Defining functions using multiple representations
\end{itemize}

\item Know, select and apply appropriate computing knowledge and problem-solving techniques. Reason about computation and systems. Use mathematical techniques to solve problems. Determine appropriate conceptual tools to apply to new situations. Know when tools do not apply and try different approaches. Critically analyze and evaluate candidate solutions.
\begin{itemize}
    \item Using the definitions of the div and mod operators on integers
\end{itemize}

\end{itemize}

\subsubsection*{TODO:}
\begin{list}
   {\itemsep2pt}
   \item \#FinAid Assignment on Canvas (complete as soon as possible) 
   \item Review quiz based on Week 0 and 1 class material (due Monday)
   \item Homework assignment 1 (due Tuesday)
   \item Review quiz based on Week 2 class material (due Monday next week)
\end{list}

\newpage




\section*{Week 2 Monday: Representing numbers}


Modeling uses data-types that are encoded in a computer.
The details of the encoding impact the efficiency of algorithms
we use to understand the systems we are modeling and the 
impacts of these algorithms on the people using the systems.
Case study: how to encode numbers?

\phantom{
Positional representation with familiar (decimal) number encodings
\vspace{30pt}
}
\vfill 

{\bf Definition} For $b$ an integer greater than $1$ and $n$ a positive integer, 
the {\bf base $b$ expansion of $n$}  is
\[
(a_{k-1} \cdots a_1 a_0)_b
\]
where $k$ is a positive integer, $a_0, a_1, \ldots, a_{k-1}$ 
are (symbols for) nonnegative integers less than $b$, $a_{k-1} \neq  0$, and
\[
n =  \sum_{i=0}^{k-1} a_{i} b^{i}
\]

Notice: {\it The base $b$ expansion of a positive integer $n$ is a string over the alphabet 
$\{x \in \mathbb{N} \st x < b\}$
whose leftmost character is nonzero.}

\begin{center}
\begin{tabular}{|c|c|}
\hline
Base $b$ & Collection of possible coefficients in base $b$ expansion of  a positive integer \\
\hline
& \\
Binary ($b=2$) & $\{0,1\}$ \\
\hline
& \\
Ternary ($b=3$) & $\{0,1, 2\}$ \\
\hline
& \\
Octal ($b=8$) & $\{0,1, 2, 3, 4, 5, 6, 7\}$\\
\hline
& \\
Decimal ($b=10$) & $\{0,1, 2, 3, 4, 5, 6, 7, 8, 9\}$\\
\hline
& \\
Hexadecimal ($b=16$) &  $\{0,1, 2, 3, 4, 5, 6, 7, 8, 9, A, B, C, D, E, F\}$\\
& letter coefficient symbols represent numerical values $(A)_{16} = (10)_{10}$\\
&$(B)_{16} = (11)_{10} ~~(C)_{16} = (12)_{10} ~~
 (D)_{16} = (13)_{10} ~~ (E)_{16} = (14)_{10} ~~ (F)_{16} = (15)_{10} $\\
\hline
\end{tabular}
\end{center}

 \vfill
\newpage




{\it Examples}:

$(1401)_{2}$

\vfill

$(1401)_{10}$

\vfill
\vfill
\vfill


$(1401)_{16}$

\vfill
\vfill
\vfill
 

{\bf Definition} For $b$ an integer greater than $1$, $w$ a positive integer, 
and $n$ a nonnegative integer
$\underline{\phantom{\hspace{1in}}}$, ~
the {\bf base $b$ fixed-width $w$ expansion of $n$}  is
\[
(a_{w-1} \cdots a_1 a_0)_{b,w}
\]
where  $a_0, a_1, \ldots, a_{w-1}$ are nonnegative integers less than $b$ and
\[
n =  \sum_{i=0}^{w-1} a_{i} b^{i}
\]
 

\begin{center}
    \begin{tabular}{|c|c|c|c|c|}
    \hline
    Decimal &  Binary  & Binary fixed-width $10$& Binary fixed-width $7$ & Binary fixed-width $4$\\
    $b=10$ & $b=2$ & $b=2$, $w =  10$& $b=2$, $w =  7$& $b=2$, $w =  4$ \\
    \hline 
    &&&&  \\
    $(20)_{10}$&\phantom{$(10100)_{2}$\qquad\qquad}&&  &\\
    &&&&  \\
\hline
    \end{tabular}
    \end{center}
 \vfill
\newpage




{\bf Definition} For $b$ an integer greater than $1$, $w$ a positive integer, 
$w'$ a positive  integer, and $x$ a real number the {\bf base $b$ fixed-width 
expansion of $x$ with integer part width $w$  and fractional part width $w'$} is
$(a_{w-1} \cdots a_1 a_0 .  c_{1} \cdots c_{w'})_{b,w,w'}$
where  $a_0, a_1, \ldots, a_{w-1}, c_1, \ldots, c_{w'}$ are nonnegative integers less than $b$ and
$$x \geq \sum_{i=0}^{w-1} a_{i} b^{i} + \sum_{j=1}^{w'}  c_{j} b^{-j} \hfill
\textrm{\qquad and \qquad}
\hfill x < \sum_{i=0}^{w-1} a_{i} b^{i} + \sum_{j=1}^{w'} c_{j} b^{-j} + b^{-w'}$$
 
\begin{center}
\begin{tabular}{|c|p{5in}|}
\hline
& \\
$3.75$  in fixed-width binary,& \\
integer part width $2$,&\\
 fractional part width $8$ & \\
& \\
& \\
& \\
& \\
\hline
& \\
$0.1$  in fixed-width binary, & \\
integer part width $2$, &\\
 fractional part width $8$ & \\
 & \\
 & \\
 & \\
 & \\
 \hline
\end{tabular}
\end{center}

\vfill

\includegraphics[width=2in]{resources/images/ArithmeticDemo.png}

Note: Java uses floating point, not fixed width representation, but similar rounding errors appear in both.
 \vfill


{\bf Representing $0$}: Which of the representations so far allow us to represent $0$?

\vfill


\begin{comment}
So far, we have representations for
positive and negative integers. What about $0$?

\begin{tabular}{|cc|p{3.4in}|p{3.7in}|}
   \hline
   & & To  represent a {\bf non-negative} integer $n$ & To represent a {\bf non-positive} integer $-n$\\
   \hline
   && &  \\
   &\parbox[t]{2mm}{\multirow{4}{*}{\rotatebox[origin=c]{90}{Sign-magnitude}}} &
   $[ 0a_{w-2} \cdots a_0]_{s,w}$, where $n =  (a_{w-2} \cdots a_0)_{2,w-1}$& 
   $[1a_{w-2} \cdots a_0]_{s,w}$
   , where $n =  (a_{w-2} \cdots a_0)_{2,w-1}$\\
   && & \\
   && Example $n=0$, $w=7$:  & Example $-n=0$, $w=7$: \\
   && & \\
   && & \\
   && & \\
   && & \\
   && & \\
   && & \\
   && & \\
   && & \\
   && & \\
   && & \\
\hline
   &&  &  \\
   &\parbox[t]{2mm}{\multirow{4}{*}{\rotatebox[origin=c]{90}{2s complement}}} &
   $[0a_{w-2} \cdots a_0]_{2c,w}$, where $n =  (a_{w-2} \cdots a_0)_{2,w-1}$& $[1a_{w-2} \cdots a_0]_{2c,w}$, where $2^{w-1} - n =  (a_{w-2} \cdots a_0)_{2,w-1}$\\
   && & \\
   && Example $n=0$, $w=7$:  & Example $-n=0$, $w=7$: \\
   && & \\
   && & \\
   && & \\
   && & \\
   && & \\
   && & \\
   && & \\
   && & \\
   && & \\
   && & \\
\hline
\end{tabular}
\end{comment} 
\newpage
\section*{Week 2 Wednesday: Algorithms for numbers}


Find and fix any and all mistakes with the following:
\begin{itemize}
\item[(a)] $(1)_2 = (1)_8$
\item[(b)] $(142)_{10} = (142)_{16}$
\item[(c)] $(20)_{10} = (10100)_2$
\item[(d)] $(35)_8 = (1D)_{16}$
\end{itemize} 

\fbox{\parbox{\textwidth}{{\bf New!} An algorithm is a finite sequence of precise instructions for solving a problem.
\hfill
}}

Algorithms can be expressed in English or in more formalized descriptions like pseudocode or fully executable programs.


Sometimes, we can define algorithms whose output matches the 
rule for a function we already care about. Consider the (integer) logarithm function
\[
logb  : \{b \in \mathbb{Z} \mid b >1 \}  \times \mathbb{Z}^+ ~~\to~~ \mathbb{N}
\]
defined by 
\[
logb (~ (b,n)~) =  \text{greatest integer } y \text{ so that } b^y  \text{ is less than or equal to } n 
\]



\begin{algorithm}[caption={Calculating integer part of base $b$ logarithm}]
   procedure $logb$($b$,$n$: positive integers with $b > 1$)
   $i$ := $0$
   while $n$ > $b-1$
     $i$ := $i + 1$
     $n$ := $n$ div $b$
   return $i$ $\{ i$ holds the integer part of the base $b$ logarithm of $n\}$
\end{algorithm} 
Trace this algorithm with inputs $b=3$ and $n=17$

\vfill
  \begin{tabular}{l|c|c|c|c}
  & $b$ & $n$ & $i$  & $n > b-1$?\\
  \hline 
  Initial value & ~$3$~ & $17$ & \phantom{$0$} & \phantom{~Yes~}\\
  After 1 iteration  & \phantom{$3$} & \phantom{$5$} & \phantom{$1$} & \phantom{T}\\
  After 2 iterations & \phantom{$3$} & \phantom{$1$} & \phantom{$2$} & \phantom{F}\\
  After 3 iterations &  &  & & \\
  &&&&\\
  \end{tabular}

\vfill

\vfill

Compare: does the output match the rule for the (integer) logarithm function?


\newpage 

{\bf Two algorithms for constructing base $b$ expansion from decimal representation}

{\bf Most significant first}: Start with left-most coefficient of expansion (highest value)

{\it Informally}: Build up to the value we need to represent in ``greedy'' approach, using 
units determined by base.

\vfill




\begin{algorithm}[caption={Calculating base $b$ expansion, from left}]
procedure $\textit{basebmost}$($n, b$: positive integers with $b > 1$)
$v$ := $n$
$k$ := $1 + $ output of $logb$ algorithm with inputs $b$ and $n$
for $i$ := $1$ to $k$
  $a_{k-i}$ := $0$
  while $v \geq b^{k-i}$
    $a_{k-i}$ := $a_{k-i} + 1$
    $v$ := $v -  b^{k-i}$
return $(a_{k-1}, \ldots, a_0) \{(a_{k-1} \ldots a_0)_b~\textrm{ is the base } b \textrm{ expansion of } n \}$
\end{algorithm}
 


\vfill
\vfill

\newpage
{\bf Least significant first}: Start with right-most coefficient of expansion (lowest value)

Idea: {\tiny(when $k > 1$)} 
    \begin{align*}
      n &= a_{k-1} b^{k-1} + \cdots + a_1 b + a_0 \\
        &= b ( a_{k-1} b^{k-2} + \cdots + a_1) + a_0\end{align*}
    so $a_0 = n \textbf{ mod } b$ and $a_{k-1} b^{k-2} + \cdots + a_1 = n \textbf{ div } b$.





\begin{algorithm}[caption={Calculating base $b$ expansion, from left}]
procedure $\textit{basebmost}$($n, b$: positive integers with $b > 1$)
$v$ := $n$
$k$ := $1 + $ output of $logb$ algorithm with inputs $b$ and $n$
for $i$ := $1$ to $k$
  $a_{k-i}$ := $0$
  while $v \geq b^{k-i}$
    $a_{k-i}$ := $a_{k-i} + 1$
    $v$ := $v -  b^{k-i}$
return $(a_{k-1}, \ldots, a_0) \{(a_{k-1} \ldots a_0)_b~\textrm{ is the base } b \textrm{ expansion of } n \}$
\end{algorithm}
 



\vfill
\vfill
\newpage 



Practice: write an algorithm for converting from base $b_1$ expansion to base $b_2$ expansion:

\phantom{
Earlier, we saw (two different) algorithms for, given 
a target base $b$, converting from decimal to base $b$ expansions. 
We will use either one of these as a subroutine in this algorithm.\\
Given a base expansion in base $b_1$:\\
Step 1: Use the definition of base expansion to calculate the value of
    this number (in decimal).\\
Step 2: Use the Least Significant First algorithm to write this value in 
    base $b_2$ and output the result.
}
\vspace{200pt} 

\newpage
\section*{Week 2 Friday: Numbers to circuits}


{\bf Representing negative integers in binary}: Fix a positive integer  width for the representation  $w$, $w >1$.

\begin{tabular}{|cc|p{3.4in}|p{3.7in}|}
\hline
& & To  represent a positive integer $n$ & To represent a negative integer $-n$\\
\hline
&& &  \\
&\parbox[t]{2mm}{\multirow{4}{*}{\rotatebox[origin=c]{90}{Sign-magnitude}}} &
$[ 0a_{w-2} \cdots a_0]_{s,w}$, where $n =  (a_{w-2} \cdots a_0)_{2,w-1}$& 
$[1a_{w-2} \cdots a_0]_{s,w}$
, where $n =  (a_{w-2} \cdots a_0)_{2,w-1}$\\
&& & \\
&& Example $n=17$, $w=7$:  & Example $-n=-17$, $w=7$: \\
&& & \\
&& & \\
&& & \\
&& & \\
&& & \\
&& & \\
&& & \\
\hline
&&  &  \\
&\parbox[t]{2mm}{\multirow{4}{*}{\rotatebox[origin=c]{90}{2s complement}}} &
$[0a_{w-2} \cdots a_0]_{2c,w}$, where $n =  (a_{w-2} \cdots a_0)_{2,w-1}$& $[1a_{w-2} \cdots a_0]_{2c,w}$, where $2^{w-1} - n =  (a_{w-2} \cdots a_0)_{2,w-1}$\\
&& & \\
&& Example $n=17$, $w=7$:  & Example $-n=-17$, $w=7$: \\
&& & \\
&& & \\
&& & \\
&& & \\
&& & \\
&& & \\
&& & \\
\hline
\end{tabular} 

For positive integer $n$, to represent $-n$ in 
2s complement with width $w$,
\begin{itemize}
    \item Calculate $2^{w-1} - n$, convert 
    result to binary fixed-width $w-1$, pad 
    with leading $1$, or
    \item Express $-n$ as a sum of powers of $2$, 
    where the leftmost $2^{w-1}$ is negative weight, or
    \item Convert $n$ to binary fixed-width $w$, 
    flip bits, add 1 (ignore overflow)
\end{itemize}

{\it Challenge: use definitions to explain why
each of these approaches works.} \newpage


{\bf Fixed-width addition}: adding one bit at time, using the usual column-by-column and carry arithmetic, and dropping the carry from the leftmost column so the result is the same width as the summands.  {\it Does this give the right value for the sum?}


\begin{align*}
   & [2~ 6]_{10,2}\\
+ &  [8~ 4]_{10,2}\\
&\overline{\phantom{[0~0]_{10,2}}}\\
\end{align*}

\begin{comment}
\begin{multicols}{2}
\begin{align*}
   & [0~ 1~ 0~ 1]_{s,4}\\
+ &  [1~ 1~ 0~ 1]_{s,4}\\
&\overline{\phantom{[0~0~1~0]_{s,4}}}\\
\end{align*}

\begin{align*}
   & [0~ 1~ 0~ 1]_{2c,4}\\
+ &  [1~ 0~ 1~ 1]_{2c,4}\\
&\overline{\phantom{[0~0~0~0]_{2c,4}}}\\
\end{align*}

\end{multicols}
\end{comment}
\vfill

\begin{multicols}{3}
   \begin{align*}
      & (1~ 1~ 0~ 1~ 0~ 0)_{2,6}\\
   + & (0~ 0~ 0~ 1~ 0~ 1)_{2,6}\\
   &\overline{\phantom{(1~1~1~0~0~1)_{2,6}}}\\
   \end{align*}
   
   \begin{align*}
      & [1~ 1~ 0~ 1~ 0~ 0]_{s,6}\\
   + & [0~ 0~ 0~ 1~ 0~ 1]_{s,6}\\
   &\overline{\phantom{(1~1~1~0~0~1)_2}}\\
   \end{align*}
   
   \begin{align*}
      & [1~ 1~ 0~ 1~ 0~ 0]_{2c,6}\\
   + & [0~ 0~ 0~ 1~ 0~ 1]_{2c,6}\\
   &\overline{\phantom{(1~1~1~0~0~1)_2}}\\
   \end{align*}
\end{multicols}
\vfill

Can fixed-width addition be used to calculate $17 + (-17) = 0$?
Hint: use 5 bits and 2s complement.

\vfill


\newpage \vfill
\vfill


In a {\bf combinatorial circuit} (also known as
a {\bf logic circuit}), we have {\bf logic gates} 
connected
by {\bf wires}. The inputs to the circuits are the 
values set on the input wires: possible
values are 0 (low) or 1 (high). The values
flow along the wires from left to right.
A wire may be split into two or more wires, 
indicated with a filled-in circle (representing
solder). Values stay the same along a wire. When 
one or more wires flow into a gate, the output 
value of that gate is computed
from the input values based on the gate's definition
table. Outputs of gates may become inputs to other
gates.  


\begin{multicols}{2}
\begin{center}\begin{tabular}{cc|c}
Inputs &  & Output \\
$x$ & $y$ & $x \text{ AND } y$  \\
\hline
$1$ & $1$ & $1$\\
$1$ & $0$ & $0$\\
$0$ & $1$ & $0$\\
$0$ & $0$ & $0$\\
\end{tabular}\end{center}
\columnbreak
\begin{center}\includegraphics[height=0.6in]{resources/images/xANDy.png} \end{center}
\end{multicols}

\vfill

\begin{multicols}{2}
\begin{center}\begin{tabular}{cc|c}
Inputs &  & Output \\
$x$ & $y$ & $x \text{ XOR } y$  \\
\hline
$1$ & $1$ & $0$\\
$1$ & $0$ & $1$\\
$0$ & $1$ & $1$\\
$0$ & $0$ & $0$\\
\end{tabular}\end{center}
\columnbreak
\begin{center}\includegraphics[height=0.4in]{resources/images/xXORy.png} \end{center}
\end{multicols}

\vfill

\begin{multicols}{2}
\begin{center}\begin{tabular}{c|c}
Input  & Output \\
$x$ & $\text{NOT } x$  \\
\hline
$1$ & $0$\\
$0$ & $1$\\
\end{tabular}\end{center}
\columnbreak
\begin{center}\includegraphics[height=0.5in]{resources/images/NOTx.png} \end{center}
\end{multicols}

%
 \vfill


{\bf Example digital circuit}: 

\begin{multicols}{2}
\begin{center}
   \includegraphics[width=1.2in]{resources/images/circuitEx.png} 
\end{center}
\columnbreak
Output when $x=1, y=0, z=0, w = 1$ is \underline{\phantom{$~~~0~~~$}}
Output when $x=1, y=1, z=1, w = 1$ is \underline{\phantom{$~~~0~~~$}}
Output when $x=0, y=0, z=0, w = 1$ is \underline{\phantom{$~~~0~~~$}}
\phantom{Output when $x=0, y=0, z=0, w = 0$ is \underline{\phantom{$~~~0~~~$}}}
\end{multicols}



Draw a logic circuit with inputs $x$ and $y$ whose output  is always $0$.  {\it  Can you use exactly 1 gate?}


\vspace{40pt} \vfill
\newpage

\begin{comment}

\newpage
\subsection*{Review Quiz}
\begin{enumerate}
\item Functions and algorithms
\begin{enumerate}
    \item {

What is a recursive definition of the set
$\mathbb{Z}^+ \times \mathbb{N}$ that is the domain of the 
function ``$b$ to the power of $i$''?
For convenience, we'll refer to this set as $X$ 
in the options below.

\begin{enumerate}[labelindent=-10pt, leftmargin=-10pt]
    \item[] Basis step: $(1,0) \in X$. Recursive step: If $(m,n) \in X$ then $(m+1, n+1) \in X$ too.
    \item[] Basis step: $(1,0) \in X$. Recursive step: If $(m,m-1) \in X$ then $(m+1, m) \in X$ too.
    \item[] Basis step: $(b,0) \in X$ for each $b \in \mathbb{Z}^+$. Recursive step: If $(m,n) \in X$ then $(m+1, n+1) \in X$ too.
    \item[] Basis step: $(b,0) \in X$ for each $b \in \mathbb{Z}^+$. Recursive step: If $(m,n) \in X$ then $(m, n+1) \in X$ too.
    \item[] None of the above.
\end{enumerate} }
    \item {

When running the algorithm $logb$ for calculating the 
integer part of base $b$ logarithm with inputs $b=4$ and $n=25$, 
which of the following calculations are helpful? Select all and only 
the calculations that are both relevant to the algorithm trace {\bf and}
are correct.

\begin{enumerate}
    \item[] $25 \textbf{ div } 4 = 5$
    \item[] $25 \textbf{ div } 4 = 6$
    \item[] $25 \textbf{ div } 4 = 1$
    \item[] $4 \textbf{ div } 25 = 0$
    \item[] $4 \textbf{ div } 25 = 5$
    \item[] $4 \textbf{ div } 25 = 1$
    \item[] $6 \textbf{ div } 4= 1$
    \item[] $5 \textbf{ div } 4= 1$
    \item[] $4 \textbf{ div } 4= 1$
\end{enumerate} }
\end{enumerate}
\item Base expansions
\begin{enumerate}
    \item {

Give the value (using usual mathematical conventions) of 
each of the following base expansions.
\begin{enumerate}
    \item $(10)_{2}$
    \item $(10)_{4}$
    \item $(17)_{16}$
    \item $(211)_{3}$
    \item $(3)_{8}$
\end{enumerate} }
    \item {

Recall the definitions from class for number representations for {\bf base $b$ expansion of $n$},
{\bf  base $b$ fixed-width $w$ expansion of $n$}, and {\bf base $b$ fixed-width expansion of $x$ 
with integer part width $w$ and fractional part width $w'$}.

For example, the base $2$ (binary) expansion of $4$ is 
$\qquad
(100)_2 \qquad$
and the base $2$ (binary) fixed-width $8$ expansion of $4$ is
$\qquad
(00000100)_{2,8} \qquad$
and the base $2$ (binary) fixed-width expansion of $4$ with integer part width $3$ and fractional
part width $2$ of $4$ is
$\qquad
(100.00)_{2,3,2} \qquad$

Compute the listed expansions.  Enter your number using the notation for base 
expansions with parentheses but without subscripts. For example, 
if your answer were $(100)_{2,3}$
you would type \texttt{(100)2,3} into Gradescope.

\begin{enumerate}
\item Give the binary (base $2$) expansion of the number whose octal (base $8$) expansion is
\[
(371)_8
\]
\item Give the decimal (base $10$) expansion of the number whose octal (base $8$) expansion is
\[
(371)_8
\]
\item Give the octal (base $8$) fixed-width $3$ expansion of $(9)_{10}$ .
\item Give the ternary (base $3$) fixed-width $8$ expansion of $(9)_{10}$ .
\item Give the hexadecimal (base $16$) fixed-width $6$ expansion of
$(16711935)_{10}$ .\footnote{This matches a frequent debugging task --
sometimes a program will show a number formatted as a base $10$
integer that is much better understood with another representation.}
\item Give the hexadecimal (base $16$) fixed-width $4$ expansion of
$$(1011~ 1010 ~ 1001~ 0000 )_2$$
Note: the spaces between each group of 4 bits above are for your convenience only.  How
might they help your calculations?
\item Give the binary fixed width expansion of $0.125$ with integer part width $2$ and 
fractional part width $4$.
\item Give the binary fixed width expansion of $1$ with integer part width $2$ and 
fractional part width $3$.
\end{enumerate} }
    \item {

Select all and only the correct choices below.
\begin{enumerate}
\item Suppose you were told that the positive integer $n_1$ has the property that $n_1 \textbf{ div } 2 = 0$. Which of the following can you conclude?
\begin{enumerate}
\item $n_1$ has a binary (base $2$) expansion
\item $n_1$ has a ternary (base $3$) expansion
\item $n_1$ has a hexadecimal (base $16$) expansion
\item $n_1$ has a base $2$ fixed-width $1$ expansion
\item $n_1$ has a base $2$ fixed-width $20$ expansion
\end{enumerate}
\item Suppose you were told that the positive integer $n_2$ has the property that $n_2 \textbf{ mod } 4 = 0$. Which of the following can you conclude?
\begin{enumerate}
\item the leftmost symbol in the binary (base $2$) expansion of $n_2$ is $1$
\item the leftmost symbol in the base $4$ expansion of $n_2$ is $1$
\item the rightmost symbol in the base $4$ expansion of $n_2$ is $0$
\item the rightmost symbol in the octal (base $8$) expansion of $n_2$ is $0$
\end{enumerate}
\end{enumerate} }
    \item {

Recall the definitions of signed integer representations from class: 
sign-magnitude and 2s complement.

\begin{enumerate}
    \item Give the 2s complement width 6 representation of the number 
    represented in binary fixed-width 5
    representation as $(00101)_{2,5}$. 
    \item Give the 2s complement width 6 representation of the number 
    represented in binary fixed-width 5
    representation as $(10101)_{2,5}$. 
    \item Give the 2s complement width 4 representation of the number 
    represented in sign-magnitude
    width 4 as $[1111]_{s,4}$.
    \item Give the sign magnitude width 4 representation of the number 
    represented in 2s complement
    width 4 as $[1111]_{2c,4}$.
    \item Give the sign magnitude width 6 representation of the number 
    represented in sign magnitude width 4 as $[1111]_{s,4}$.
    \item Give the 2s complement width 6 representation of the number 
    represented in 2s complement width 4 as $[1111]_{2c,4}$.
\end{enumerate} }
\end{enumerate}
\item Multiple representations
{\input{../activity-snippets/color-hexcolor-intro-question.tex}}
\end{enumerate}
\end{comment}

\end{document}
